% ============================================================
%  Chapter fragment -- included by main.tex
%  (not a standalone document; compile main.tex to build the PDF)
% ============================================================
\begin{multicols*}{3}
\raggedcolumns

\sheettitle{01b -- Securing Information Systems}{SWS2 -- Overview / ZHAW / Tellenbach}

% ==================== CHAPTER 1: SECURING IS ====================
\section{Securing Information Systems}

\subsection{Information System (ISO 27000)}
Set of applications, services, IT assets \& info-handling components. Involves:
\begin{itemize}
\item \kw{Methods \& procedures}
\item \kw{Data}
\item \kw{Technology} (software/hardware)
\item \kw{People}
\end{itemize}

\subsection{Two fundamental questions}
\begin{itemize}
\item How to \kw{manage security} so the system stays secure?
\item How do we \kw{know} we are secure? (measuring)
\end{itemize}

\subsection{ISMS (Info Sec Mgmt System)}
\begin{itemize}
\item \kw{Systematic approach} to managing info so it stays secure.
\item Covers people, processes, IT via a \kw{risk management process}.
\item Risk kept at an \kw{acceptable level} through a coherent set of \kw{security controls}.
\item Examples: \term{ISO 27000-series, NIST RMF, BSI IT-Grundschutz, CIS Controls}.
\end{itemize}

\subsection{Security Controls}
Safeguards/countermeasures to \kw{avoid, detect, counteract, minimize} risks to assets.
Characterized by \kw{attributes}:
\begin{itemize}
\item \kw{Info-sec property}: Confidentiality, Integrity, Availability (CIA)
\item \kw{Category}: people, physical, technology, organizational
\item \kw{Function}: identify, protect, detect, respond, recover
\end{itemize}

\subsection{Type of a Control \hl{(!)}}
\begin{itemize}
\item \kw{Preventive}: prevent incident (auth, firewall)
\item \kw{Detective}: identify incident in progress (alarm, IDS)
\item \kw{Corrective}: limit damage / recover (backup, redundancy)
\end{itemize}

\subsection{Controls in ISMS -- Abstraction}
\begin{itemize}
\item ISMS operate at \kw{high abstraction} (apply to all orgs/sectors).
\item Different granularity; sector-specific control sets exist.
\item Little guidance on \kw{technical} implementation.
\end{itemize}

\subsection{Example: ISO 27000 ("27k")}
\begin{itemize}
\item \kw{ISO 27001:2022}: 93 controls in Annex A (For Example): Network Security: "Networks and network devices should be secured, managed and controlled to protect information in systems and applications." So it says what to do, but not how to do it.
\item \kw{ISO 27002:2022}: adds \kw{implementation guidance} -- control objectives \& best-practice controls.
\item Not all controls relevant to every org; industry extensions exist (e.g. telco ISO 27011).
\end{itemize}
ISO 27002 control documents have: \term{Control, Purpose, Guidance} + attribute table.

\subsection{Ex: Control "Protection against malware"}
ISO 27002:2022 tags each control with 5 \kw{attributes}; possible values per attribute below, \hl{highlighted = the ones that apply} for this control:
\begin{itemize}
\item \kw{Control type}: Preventive / \hl{Detective} / Corrective
\item \kw{Info security properties}: Confidentiality / \hl{Integrity} / Availability
\item \kw{Cybersecurity concepts}: Protect / \hl{Detect}
\item \kw{Operational capabilities}: System\_and\_network\_security / Information\_protection
\item \kw{Security domains}: Protection / \hl{Defence}
\end{itemize}

\subsection{Implementing an ISMS}
\begin{itemize}
\item Task for \kw{management} -- usually the \kw{CISO} (talks to board in business/risk terms).
\item Gives the \kw{big picture}; standards act as a \kw{checklist}.
\item \kw{Certification}: checks whether company implements a standard.
\item Complex, costly; \hl{no guidance to efficiently prioritize} vs most common threats $\Rightarrow$ \kw{CIS Controls}.
\end{itemize}

% ==================== CHAPTER: CIS CONTROLS ====================
\section{CIS Critical Security Controls}

\subsection{Overview}
Is a list of the most important security measures to defend against real cyber-attacks

In comparison to ISO 27000, ISO is high-level and says "what" to do, but not "how" to do it. CIS Controls are more specific and actionable, with detailed implementation guidance.
\begin{itemize}
\item Best-practice guidelines, started \kw{2008} (US defense breaches).
\item \kw{18 controls} for current cyber attacks; by experts.
\end{itemize}
\textbf{Guiding principles:}
\begin{itemize}
\item Only \kw{critical} controls, based on \kw{attacker behavior} \& data.
\item \kw{Prioritize as a community}; immediate action $>$ completeness.
\item Provide info on how to implement \& \kw{automate}.
\end{itemize}

\subsection{The 18 Controls}
{\scriptsize
01 Inventory Enterprise Assets \textbullet{} 02 Inventory Software Assets \textbullet{} 03 Data Protection \textbullet{} 04 Secure Config \textbullet{} 05 Account Mgmt \textbullet{} 06 Access Control Mgmt \textbullet{} 07 Continuous Vuln Mgmt \textbullet{} 08 Audit Log Mgmt \textbullet{} 09 Email \& Web Browser \textbullet{} 10 Malware Defenses \textbullet{} 11 Data Recovery \textbullet{} 12 Network Infra Mgmt \textbullet{} 13 Network Monitoring \& Defense \textbullet{} 14 Security Awareness \& Training \textbullet{} 15 Service Provider Mgmt \textbullet{} 16 Application Software Sec \textbullet{} 17 Incident Response Mgmt \textbullet{} 18 Penetration Testing
}

\subsection{Safeguards}
\kw{Safeguards} = concrete things to implement (called sub-controls before v8). Each control has N safeguards split across IG1/IG2/IG3.

\subsection{CIS Control content}
\begin{itemize}
\item \kw{Overview}: purpose \& defensive use
\item \kw{Why important}: how attackers exploit absence
\item \kw{Processes \& tools}: technical impl./automation
\item \kw{Safeguards}: table of specific measures
\end{itemize}

\subsection{Prioritization \hl{(!)}}
\begin{itemize}
\item \kw{Until v7.1}: controls 1--6 = "cyber hygiene" first. Criticized as too simplistic.
\item \kw{Since v8}: prioritization via \kw{Implementation Groups IG1--IG3}. Control order still has value (e.g. 1\&2 protect vs malware better than 10).
\item In practice: baseline + \hl{risk assessment} (cost-benefit per environment).
\end{itemize}

\subsection{Implementation Groups (IG)}
\begin{itemize}
\item \kw{IG1 -- Basic Cyber Hygiene}: limited expertise, COTS, small biz/home; thwarts \kw{general, non-targeted} attacks. Stop here if SME, low-sensitivity data, focus on operations.
\item \kw{IG2} (incl IG1): copes with operational complexity; enterprise-grade tech \& expertise; multiple depts, compliance, sensitive data, public trust.
\item \kw{IG3} (incl IG1+IG2): security experts; vs \kw{targeted attacks}, reduce impact of \kw{zero-days}; regulated, availability+CIA at heart of business.
\end{itemize}

\subsection{Safeguards per IG -- "x/N" notation \hl{(!)}}
Each control has N safeguards; IGs are \kw{cumulative} (IG2$\supseteq$IG1, IG3$\supseteq$IG2). "\kw{x/N}" = how many of the N safeguards apply at that level.
\textit{Ex: Control 09 Email \& Web Browser (7 safeguards)} $\to$ IG1 \hl{2/7}, IG2 6/7, IG3 7/7:
\begin{itemize}
\item \kw{IG1 (2/7)}: only supported browsers/mail clients + DNS filtering (cheap, high-impact).
\item \kw{IG2 (6/7)}: + URL filters, restrict extensions, DMARC, block file types.
\item \kw{IG3 (7/7)}: + email-server anti-malware.
\end{itemize}
\kw{IG1 = 0/N} (e.g. 13 Net Monitoring 0/11, 18 PenTest 0/5) $\Rightarrow$ basic-hygiene org does \kw{nothing} in that control yet.

\subsection{Selected Controls (detail)}
\begin{itemize}
\item \kw{CSC 1 -- Enterprise Assets}: active+passive device discovery, asset inventory DB, NLA (802.1x), PKI, analytics. \term{Why both active+passive?} active=fast/triggers, passive=continuous/finds hidden \& unauth devices.
\item \kw{CSC 2 -- Software Assets}: software whitelisting, inventory tool, NAC to \kw{enforce}. Very effective vs malware (blocks unauthorized SW).
\item \kw{CSC 7 -- Continuous Vuln Mgmt}: SCAP vuln scanner + patch mgmt + analytics.
\end{itemize}

\subsection{Summary \& Criticism}
\begin{itemize}
\item Most companies implement only a \kw{subset}.
\item Protecting IS is complex: needs checklists, resources, risk mgmt.
\item Critique: paper "A Critical View on CIS Controls" (S. Gros) -- some valid, some academic.
\end{itemize}

% ==================== CHAPTER: MEASURING SECURITY ====================
\section{Measuring Security of IS}

\subsection{How secure is your org?}
Factors (incomplete): \kw{people, physical security, vulnerabilities (known+unknown), controls, configuration \& architecture, processes \& policies}.
\hl{Measuring security scientifically/holistically is (most likely) not possible.} Same for measuring \kw{risk}.

\subsection{Approaches to measure}
\begin{itemize}
\item \kw{Security}: \kw{Audit} (systematic eval vs established criteria), \kw{penetration test}, measures of ISMS performance \& effectiveness.
\item \kw{Risk}: \kw{Risk = Likelihood $\times$ Impact}; estimating likelihood is very hard.
\end{itemize}

\subsection{Why measure at all?}
\begin{itemize}
\item Identify controls \kw{implemented incorrectly / ineffective / missing}.
\item \kw{Quantify progress} (e.g. time to fix vuln, phishing resilience).
\item \kw{Compliance} evidence (ISO 27001 requires eval of performance \& effectiveness).
\item Support \kw{risk-informed decisions}.
\end{itemize}

\subsection{Measure frameworks}
\kw{NIST SP 800-55} \& \kw{ISO 27004} provide measure templates (Goal, Measure, Formula, Target, Frequency, Responsible parties, Data source, Reporting format). E.g. \% of high vulns mitigated within defined time; trend of detected-but-not-blocked malware attacks.

\subsection{Recommendations}
\begin{itemize}
\item Measuring whole infrastructure safety hardly possible, but we must measure to \kw{decide} $\Rightarrow$ aim: \kw{reduce uncertainty}.
\item \kw{Question} what/how is measured \& what conclusions (e.g. "2$\times$ failed SSH logins" $\neq$ 2$\times$ exposure).
\item Use \kw{comparative measurements} (IDS subnet $>$ no IDS; static analysis $\to$ more secure).
\item Measure \kw{progress \& maturity} with best-practice measures.
\end{itemize}

\begin{defbox}
\textbf{Take-away:} Securing information systems is hard. Measuring security is hard.
\end{defbox}

\end{multicols*}
